\hwpchapter{3}{영향요인 및 오류분석}{15}

개발 구간(f1--f4, 53일, 5,016슬롯, 피크 평가일 51일)에서 C-BAT가 잘 맞는 조건과 틀리는 조건을 구분하였다. 슬롯 오차는 예측 중앙값으로, 일 피크 오차는 사후 경로마다 구한 하루 최댓값(이하 경로 최댓값)의 중앙값으로 계산하였으며, C-BAT 수치는 seed 0--2의 평균이다. 고피크일은 일 피크가 해당 학습 기간 C90을 넘은 13일이다. 지표의 정의는 2.2절을 따른다. 9월 시험 구간은 고피크일이 없고 비가동일이 2일뿐이어서 조건별 분석에는 쓰지 않고 보조로만 언급하였다.

\hwpsection{3.1 생산조건별 예측오차}

슬롯을 그날의 가동유형과 그 슬롯 시각의 계획 생산 여부로 일곱 조건으로 나누어 슬롯 MAE를 비교하였다(표~\ref{tab:ch3-regime}). 비생산은 해당 시각의 계획 생산량이 0인 상태이며 설비 전체의 정지를 뜻하지 않는다. C-BAT 수치는 C-BAT의 조건별 오차 산출 파일에서, 기존 BAT·M2·B0 수치는 같은 개발 구간 예측의 조건별 오차 파일에서 가져왔다. 두 파일의 조건별 관측 슬롯 수는 일곱 조건 모두 같았다(합계 5,016).

\begin{table}[htbp]
\centering
\hwptablefont
\setlength{\tabcolsep}{4pt}
\caption{생산조건별 슬롯 MAE(개발 구간 53일, kW). C-BAT 편향은 예측 중앙값에서 관측을 뺀 값의 평균이다.}\label{tab:ch3-regime}
\begin{tabular}{L{5.4cm}rrrrrrr}
\toprule
 & & & \multicolumn{2}{c}{C-BAT} & & & \\
\cmidrule(lr){4-5}
생산조건 & 관측 슬롯 & 일수 & MAE & 편향 & 기존 BAT & M2 & B0 \\
\midrule
가동유형 0 (비가동) & 1,394 & 15 & 1.73 & $+0.73$ & 1.67 & 4.41 & 34.32 \\
가동유형 1 (1--12시간)·비생산 & 378 & 6 & 7.95 & $-6.14$ & 4.06 & 3.33 & 3.75 \\
가동유형 1 (1--12시간)·생산 & 172 & 6 & 8.69 & $-5.84$ & 7.71 & 11.26 & 17.26 \\
가동유형 2 (13--19시간)·비생산 & 268 & 8 & 20.66 & $+13.67$ & 9.29 & 15.42 & 6.18 \\
가동유형 2 (13--19시간)·생산 & 500 & 8 & 14.52 & $-2.38$ & 16.15 & 17.97 & 28.25 \\
가동유형 3 (20시간 이상)·비생산 & 108 & 16 & 7.98 & $-5.09$ & 8.00 & 18.66 & 29.19 \\
가동유형 3 (20시간 이상)·생산 & 2,196 & 24 & 10.05 & $-2.92$ & 10.30 & 15.43 & 28.90 \\
\midrule
전체 & 5,016 & 53 & 8.50 & -- & 7.82 & 11.64 & 26.84 \\
\bottomrule
\end{tabular}
\end{table}
\hwpnote{관측 슬롯은 목표가 관측된 15분 슬롯 수이고, 일수는 그 조건의 슬롯을 하나 이상 가진 날 수이다. 가동유형 3·비생산은 20시간 이상 가동일 안의 생산 공백 시각으로, 16일에 걸쳐 108슬롯뿐이다. C-BAT는 seed 0--2 평균이고 기존 BAT·M2·B0는 저장된 예측 한 벌이다. 편향이 양수이면 과대예측이다.}

C-BAT의 오차는 비가동일에서 1.73 kW로 가장 작고, 가동유형 2의 비생산 슬롯에서 20.66 kW로 가장 컸다. 이 슬롯의 편향은 $+13.67$ kW로, 13--19시간 가동일의 생산 시작 전과 종료 뒤 전력을 가동 수준으로 높게 예측하였다. 이 조건에서 C-BAT는 네 모형 가운데 가장 나빴고 B0(6.18 kW)보다도 오차가 컸다. C-BAT의 기준일은 가동 여부와 날짜유형만 맞추므로(2.1절) 13--19시간 가동일이 생산 시각이 다른 날의 곡선을 빌려 올 수 있으며, 이를 과대예측의 원인 후보로 보았다. 가동유형 1에서도 C-BAT(7.95, 8.69 kW)는 기존 BAT(4.06, 7.71 kW)보다 오차가 컸고, 편향이 $-6.14$, $-5.84$ kW로 하루 전체를 낮게 예측하였다. 반면 20시간 이상 가동일에서는 생산·비생산 슬롯 모두 C-BAT가 기존 BAT와 같거나 조금 작았고 M2(15.43, 18.66 kW)와 B0보다 작았다. 전체 슬롯 MAE에서 C-BAT(8.50 kW)가 기존 BAT(7.82 kW)보다 큰 것은 대부분 가동유형 2·비생산과 가동유형 1·비생산 슬롯에서 생긴 차이이며, 가동유형 2·생산과 3·생산에서는 C-BAT가 더 작았다.

생산 시작 시각별로는 07시에 시작한 2일의 슬롯 MAE가 21.56 kW(편향 $+17.80$ kW), 08시에 시작한 5일이 13.51 kW로, 전날부터 이어 가동한 20일(10.38 kW)이나 01시에 시작한 11일(9.43 kW)보다 컸다. 07--08시에 시작한 날은 모두 가동유형 2이므로 앞의 비생산 슬롯 과대예측과 같은 현상이다. 하루 중에서는 18시(11.11 kW)와 19시(10.11 kW)의 오차가 가장 크고 21--23시(5.93--6.05 kW)가 가장 작았으며, 08--16시에는 편향이 $-1.42$ kW에서 $-3.83$ kW 사이로 생산 시간의 전력을 조금 낮게 예측하였다. 요일별로는 월요일 8일이 14.10 kW(편향 $+5.22$ kW)로 가장 컸다. 일 생산량 사분위별 슬롯 MAE는 10.07--11.69 kW로 차이가 작아, 슬롯 오차는 생산량의 크기보다 생산 시각의 배치에 따라 갈렸다.

\hwpsection{3.2 일 피크 층별 오차와 개별 날짜}

피크 평가일 51일을 비가동 14일과 가동 37일로 나누고, 가동일에 포함되는 고피크일 13일을 따로 묶어 일 피크 MAE와 편향을 비교하였다(표~\ref{tab:ch3-strata}, 그림~\ref{fig:ch3-strata}). C-BAT의 피크 MAE는 비가동일 2.50 kW, 가동일 11.36 kW, 고피크일 4.37 kW로, 고피크일 오차가 가동일 전체보다 작았다. 기존 BAT는 비가동일 피크를 평균 24.59 kW 높게, LightGBM(튜닝)은 고피크일 피크를 20.37 kW 낮게 예측하였다. C-BAT는 이 두 실패를 함께 줄였으나 가동일 피크를 평균 9.74 kW 높게 예측하였다. B0\_kind는 전체(6.67 kW)·비가동(0.93 kW)·가동(8.84 kW) 피크 MAE가 가장 작았고 고피크일에서는 7.54 kW(편향 $-4.77$ kW)로 C-BAT보다 컸다. 다만 고피크일에서 C-BAT와 B0\_kind의 차이($-3.17$ kW)의 날짜 단위 짝지은 bootstrap 95\% 신뢰구간 [$-7.04$, 0.40]은 0을 포함한다.

\begin{table}[htbp]
\centering
\hwptablefont
\setlength{\tabcolsep}{4pt}
\caption{일 피크 층별 MAE와 편향(개발 구간 피크 평가일 51일, kW). 편향은 경로 최댓값의 평균에서 관측 피크를 뺀 값이다. 고피크일 13일은 가동일 37일에 포함된다.}\label{tab:ch3-strata}
\begin{tabular}{lrrrrrrrr}
\toprule
 & \multicolumn{4}{c}{피크 MAE} & \multicolumn{4}{c}{피크 편향} \\
\cmidrule(lr){2-5}\cmidrule(lr){6-9}
모형 & 전체 & 비가동 & 가동 & 고피크 & 전체 & 비가동 & 가동 & 고피크 \\
\midrule
C-BAT & 8.93 & 2.50 & 11.36 & 4.37 & $+7.77$ & $+2.54$ & $+9.74$ & $-1.08$ \\
기존 BAT & 14.50 & 24.11 & 10.86 & 12.47 & $+8.10$ & $+24.59$ & $+1.87$ & $-12.24$ \\
B0\_kind & 6.67 & 0.93 & 8.84 & 7.54 & $+0.24$ & $+0.36$ & $+0.19$ & $-4.77$ \\
LightGBM(튜닝) & 10.06 & 2.51 & 12.92 & 18.27 & $+1.98$ & $+19.47$ & $-4.64$ & $-20.37$ \\
\bottomrule
\end{tabular}
\end{table}

\begin{figure}[htbp]
\centering
\includegraphics[width=\linewidth]{peak_strata.pdf}
\caption{일 피크 층별 피크 MAE(개발 구간). 가로축은 왼쪽부터 전체 51일(All), 비가동 14일(Nonop), 가동 37일(Op), 고피크 13일(High)이다. 범례의 BAT는 기존 BAT, LightGBM은 LightGBM(튜닝)이며 막대 위 숫자는 MAE(kW)이다.}\label{fig:ch3-strata}
\end{figure}

가동일 과대예측은 평균 곡선이 아니라 경로 최댓값에서 생겼다. 가동일 37일에서 C-BAT 예측 평균 곡선의 하루 최댓값은 관측 피크보다 평균 10.90 kW 낮았다. 경로 최댓값의 중앙값은 여기에 약 20 kW를 더해 가동일 편향(중앙값 기준)을 $+9.22$ kW로 만들었다. 이 추가분은 고피크일에는 알맞았다. 고피크일 13일의 평균 곡선 최댓값은 관측보다 23.51 kW 낮았고, 경로 최댓값을 쓰면 편향이 $-1.70$ kW로 줄었다. 그러나 고피크가 아닌 가동일 24일은 평균 곡선 최댓값의 편향이 $-4.07$ kW뿐이어서, 비슷한 폭을 더한 경로 최댓값의 편향은 $+15.14$ kW가 되었다. C-BAT는 어느 가동일에 피크가 크게 오를지 구분하지 못한 채 모든 가동일에 비슷한 상승폭을 더한 것이다. 슬롯 단위 편향이 음수인 가동유형 1과 3에서도 일 피크 편향(중앙값 기준)이 각각 $+9.86$, $+8.74$ kW로 양수인 것은 이 때문이다. C-BAT는 51일 중 41일, 가동일 37일 중 28일의 피크를 높게 예측하였고, 고피크일 13일 가운데서는 9일을 낮게 예측하였다.

피크 오차가 가장 큰 10일은 모두 고피크가 아닌 가동일이었고, 모두 16.63--31.90 kW 과대예측이었다(표~\ref{tab:ch3-worst}). 관측 피크는 173--195 kW로 각 fold의 C90(198.0--211.0 kW, 1.2절)보다 낮았지만 C-BAT는 204.19--216.19 kW를 예측하였다. 10일 중 7일은 일 슬롯 MAE가 7.38--11.32 kW로 곡선은 대체로 맞았고 피크만 높았다. 일 슬롯 MAE가 큰 8월 16일(25.90 kW)과 7월 9일(22.61 kW)은 가동유형 2이며, 8월 16일은 07시에 시작한 공휴일이다(C-BAT는 공휴일을 입력으로 쓰지 않는다, 1.1절). 10일 중 9일이 8월이었고 그중 8일이 8월 16일 이후였다. 관측 고피크는 7월에 몰리고(22일 중 11일) 8월에는 적었으므로(29일 중 2일), 7월의 높은 피크를 학습한 모형이 8월의 낮아진 피크 수준을 늦게 따라간 것으로 볼 수 있다. 다만 고피크 판정 기준인 C90이 fold마다 달라 생긴 차이가 섞여 있고 표본이 작아 확정하지 않았다.

\begin{table}[htbp]
\centering
\hwptablefont
\setlength{\tabcolsep}{4pt}
\caption{C-BAT 피크 오차 상위 10일(개발 구간, seed 0--2 평균, kW). 시작 0시는 전날부터 이어 가동한 날이다.}\label{tab:ch3-worst}
\begin{tabular}{lcccrrrrr}
\toprule
날짜 & 요일 & 가동유형 & 생산 시작 & 생산시간 & 관측 피크 & 예측 피크 & 피크 오차 & 일 슬롯 MAE \\
\midrule
8월 17일 & 화 & 3 & 0시 & 23 & 173 & 204.90 & $+31.90$ & 8.06 \\
8월 16일 & 월 & 2 & 07시 & 15 & 190 & 215.58 & $+25.58$ & 25.90 \\
7월 9일 & 금 & 2 & 01시 & 18 & 192 & 216.19 & $+24.19$ & 22.61 \\
8월 31일 & 화 & 3 & 0시 & 24 & 187 & 208.35 & $+21.35$ & 11.32 \\
8월 25일 & 수 & 3 & 01시 & 23 & 189 & 210.24 & $+21.24$ & 7.54 \\
8월 27일 & 금 & 3 & 01시 & 23 & 190 & 210.92 & $+20.92$ & 10.47 \\
8월 19일 & 목 & 3 & 0시 & 24 & 188 & 208.57 & $+20.57$ & 7.38 \\
8월 18일 & 수 & 3 & 0시 & 24 & 188 & 207.37 & $+19.37$ & 10.07 \\
8월 30일 & 월 & 2 & 08시 & 16 & 186 & 204.19 & $+18.19$ & 13.43 \\
8월 10일 & 화 & 3 & 01시 & 21 & 195 & 211.63 & $+16.63$ & 7.61 \\
\bottomrule
\end{tabular}
\end{table}
\hwpnote{예측 피크는 경로 최댓값의 중앙값, 피크 오차는 예측 피크에서 관측 피크를 뺀 값이다. 10일 모두 세 seed에서 과대예측이었다. 생산시간은 생산량이 양수인 시간 수이다.}

피크 시각 적중률(실제 피크 시각의 ±2슬롯, 즉 ±30분 이내)은 51일 전체 0.203, 가동일 0.252, 고피크일 0.385였다. 비가동일은 곡선이 평탄하여 0.071에 그쳤고, 기존 BAT의 전체 적중률은 0.255였다. 고피크일의 피크가 08시와 11시에 몰린다는 관측(1.2절)이 있으나 C-BAT는 고피크일 13일 중 5일만 ±30분 안에 맞혔다. 따라서 C-BAT의 일 피크 예측은 크기 판단에 쓰고, 특정 15분 슬롯의 부하를 제어하는 근거로는 쓰지 않았다.

\hwpsection{3.3 피크 기준 초과 경보의 미탐지·오경보}

피크 기준 초과 경보는 그날 일 피크가 C50·C75·C90을 넘을 확률이 문턱 이상일 때 내는 경보이다(기준값과 초과확률의 정의는 2.1절). 개발 구간에서는 fold마다 내부 보정일 7일로 Platt 보정을 적합하고, 같은 fold 내부에서 F1이 가장 큰 확률 문턱을 정하였다. 실제 초과일에 경보가 없으면 미탐지, 초과하지 않은 날에 경보가 나면 오경보로 셌다. 51일 중 초과일은 C50 31일, C75 21일, C90 13일이다(표~\ref{tab:ch3-alarm}).

\begin{table}[htbp]
\centering
\hwptablefont
\caption{피크 기준 초과 경보 성능(개발 구간 51일, 7일 Platt 보정, 확률 모형은 seed 0--2 평균). 매일 경보는 모든 날에 경보를 낸 기준선이다.}\label{tab:ch3-alarm}
\begin{tabular}{llrrrrr}
\toprule
기준(초과일) & 모형 & 정밀도 & 재현율 & F1 & AUC & Brier \\
\midrule
C50 (31일) & C-BAT & 1.000 & 0.882 & 0.937 & 0.997 & 0.020 \\
 & 기존 BAT & 0.967 & 0.935 & 0.951 & 0.976 & 0.020 \\
 & LightGBM(튜닝) & 0.963 & 0.839 & 0.897 & 0.977 & 0.030 \\
 & 매일 경보 & 0.608 & 1.000 & 0.756 & -- & -- \\
\midrule
C75 (21일) & C-BAT & 0.829 & 0.921 & 0.872 & 0.844 & 0.153 \\
 & 기존 BAT & 0.769 & 0.952 & 0.851 & 0.863 & 0.211 \\
 & LightGBM(튜닝) & 0.900 & 0.429 & 0.581 & 0.798 & 0.236 \\
 & 매일 경보 & 0.412 & 1.000 & 0.583 & -- & -- \\
\midrule
C90 (13일) & C-BAT & 0.520 & 0.667 & 0.584 & 0.788 & 0.201 \\
 & 기존 BAT & 0.417 & 0.769 & 0.541 & 0.800 & 0.233 \\
 & LightGBM(튜닝) & 1.000 & 0.231 & 0.375 & 0.688 & 0.253 \\
 & 매일 경보 & 0.255 & 1.000 & 0.406 & -- & -- \\
\bottomrule
\end{tabular}
\end{table}
\hwpnote{매일 경보는 확률을 내지 않으므로 AUC와 Brier를 비워 두었다. C50·C75·C90은 계약전력이나 설비 한계가 아니라 학습 기간 가동일 일 피크의 상대 분위수이다.}

가장 드문 C90에서 C-BAT의 F1은 0.584, Brier는 0.201로 세 확률 모형 가운데 가장 좋았고, AUC(0.788)는 기존 BAT(0.800)보다 조금 낮았다. 기존 BAT는 재현율이 0.769로 높았으나 정밀도가 0.417로 경보 확인 부담이 컸다. LightGBM(튜닝)은 경보를 낸 날은 모두 맞혔지만 재현율이 0.231로 고피크 13일 중 3일만 잡았고, F1(0.375)이 매일 경보(0.406)보다도 낮았다. C75에서도 C-BAT의 F1이 0.872로 가장 높았고, C50에서는 세 모형 모두 F1이 0.89 이상이었다(C-BAT 0.937, 기존 BAT 0.951).

C-BAT의 C90 미탐지(seed 평균 4.33일)는 모두 7월이고 모두 가동유형 3이었으며, 62\%가 일 생산량 상위 25\%에 속하였다. 51일 전체에서 이 세 조건의 비중이 각각 43\%, 47\%, 25\%인 것과 비교하면 한쪽으로 크게 몰렸다. 7월에 생산량이 많은 종일 가동일의 피크가 학습 기간보다 높게 올라갔고 C-BAT가 그 상승을 다 따라가지 못한 것으로 보인다. 오경보(seed 평균 8일)는 7월과 8월이 절반씩이었고, 가동유형 3이 62\%, 가동유형 2가 38\%였다. 요일로는 금요일이 38\%였고 주말 오경보는 없었다. 오경보는 3.2절의 평범한 가동일 과대예측에서 나온 것이다. 미탐지는 여름철 부하 상승을 늦게 따라가서, 오경보는 평범한 가동일의 피크를 높게 보아서 생겼다.

경보 지표는 보정 방식에 민감하였다. 7일 Platt 보정은 적합에 쓸 날이 적다. 개발 구간 C-BAT의 C90 AUC는 보정 전 경로 확률에서 0.926이었으나 보정 뒤 0.788로 떨어졌고, Brier도 0.131에서 0.201로 나빠졌다. C75도 보정 전 Brier(0.146)가 보정 뒤(0.153)보다 낮았으며, C50만 보정 뒤가 조금 나았다(AUC 0.981에서 0.997, Brier 0.021에서 0.020). 최종 실행에서는 8월 말 내부 보정일 7일(쓸 수 있는 날 5일)에 C90 초과일이 없어 보정식이 절편만 추정하였고, 9월 가동일의 C90 위험이 모두 0.0013으로 고정되어 경보가 울릴 수 없었다. 이에 9월 위험 결과를 보기 전에 규칙을 정하여, 최종 발행 확률을 보정 전 경로 확률 $P(\text{경로 최댓값} > C)$로 바꾸었다. 표~\ref{tab:ch3-alarm}은 이전 보정 방식으로 계산한 개발 구간 결과이므로 모형 간 상대 비교로만 읽는다. 보정 전 확률은 평범한 가동일 과대예측을 그대로 이어받아, C90 초과일이 없던 9월에도 C90 위험이 최대 0.30까지 나왔다.

\hwpsection{3.4 위험조건과 보완 대상}

모형과 무관한 관측으로 고피크일의 조건을 보면, 13일은 모두 13시간 이상 가동한 날이었다(1.2절). 가동유형 2는 8일 중 3일(37.50\%), 가동유형 3은 24일 중 10일(41.67\%)이 고피크였고, 가동유형 1의 5일(평균 피크 107.60 kW)과 비가동일에는 고피크가 없었다. 피크 평가일 가운데 전날부터 이어 가동한 19일은 고피크 비율이 42.11\%이고 피크가 가장 자주 나온 시각이 11시였으며, 07--08시에 시작한 7일은 3일이 고피크이고 피크가 주로 08시에 나왔다. 01시에 시작한 11일은 고피크 비율이 18.18\%로 낮았다. 금요일은 7일 중 1일(14.29\%)만 고피크였고 토·일요일에는 고피크가 없었다. 일 생산량 1사분위(평균 피크 153.60 kW)를 빼면 2--4사분위의 평균 피크는 198.56--202.67 kW로 비슷하여, 생산량이 많을수록 피크가 계속 커지지는 않았다.

같은 조건 안에서 고피크일과 평범한 가동일을 가르는 신호는 약하였다. 가동일 37일에서 C-BAT 피크 오차(경로 최댓값의 중앙값에서 관측을 뺀 값, seed 평균)와 생산 시작 시각의 선형 상관은 $r=-0.03$, 생산 첫 2시간의 생산량 비중과는 $r=-0.01$, 생산시간과는 $r=-0.01$이었다. 이 세 특성의 선형 관계로는 오차가 큰 날을 구분할 수 없었다. 특성의 조합이나 비선형 관계는 따로 검증하지 않았으므로, 이 결과만으로 생산계획에 고피크 정보가 없다고 판단하지는 않았다.

생산계획 입력 오차의 영향은 학습 입력과 채점할 실제 전력을 그대로 두고 예측일에 넣는 계획만 바꾸어 확인하였다. 생산량 오차는 평균이 1인 로그정규 배수(로그 표준편차 0.05, 0.1, 0.2, 날짜별 10회 반복)로, 시각 오차는 계획 전체를 1시간 앞당기거나 늦추어 넣었다. C-BAT는 seed 0 하나로 실행하였다(표집 6,000회, burn-in 3,000회). 오차가 없을 때 8.49 kW였던 C-BAT 슬롯 MAE는 로그 표준편차 0.2에서도 8.50 kW로 0.04\% 늘었다. 계획을 1시간 이르게 넣으면 8.78 kW($+3.35$\%), 늦게 넣으면 8.64 kW($+1.77$\%)가 되었다. 기존 BAT($+0.44$, $+0.06$ kW)와 LightGBM(튜닝)($+0.87$, $+0.69$ kW)도 생산량 오차보다 시각 오차에 더 민감하였다. C-BAT의 피크 MAE는 모든 조건에서 8.61--8.94 kW였다. 전력이 생산량의 크기보다 가동 여부와 가동 시각을 따른다는 3.1절의 결과와 맞으며, 현장에서는 생산량보다 가동 시작·종료 시각의 입력 정확도가 중요하다. 계획 취소나 설비 고장처럼 가동 여부 자체가 바뀌는 경우는 시험하지 않았다.
\hwpnote{기존 BAT와 LightGBM(튜닝)의 입력 오차 결과는 같은 코드와 절차로 앞서 실행한 결과(표집 2,000회)를 옮겼다. 괄호 안 값은 오차가 없을 때 대비 슬롯 MAE 증가분이며, 앞이 1시간 이르게, 뒤가 1시간 늦게 넣은 경우이다.}

C-BAT의 구성을 바꾸어 보는 확장 실험 11가지를 사전 등록하여 개발 구간에서만 평가하였고, 채택된 것은 없다(표~\ref{tab:ch3-ext}). 각 실험은 결과를 보기 전에 채택 조건을 문서로 고정하였다. 조건은 실험마다 조금씩 다르지만 공통으로 C-BAT의 전체 피크 MAE(8.93 kW)와 고피크일 피크 MAE(4.37 kW)를 지키면서 곡선이나 가동일 편향을 개선하도록 요구하였다. 9월 시험 구간은 쓰지 않았다.

\begin{table}[htbp]
\centering
\hwptablefont
\setlength{\tabcolsep}{4pt}
\caption{사전 등록 확장 실험(개발 구간 f1--f4, seed 0--2 평균, kW). 앙상블은 C-BAT seed 0과 비교하였고 fold 하나씩 뺀 가중치로 평가하였으며, 일 피크는 C-BAT와 같다.}\label{tab:ch3-ext}
\begin{tabular}{L{1.4cm}L{4.1cm}rrrL{4.1cm}}
\toprule
분류 & 변형 & 슬롯 MAE & 피크 MAE & 고피크 MAE & 미충족 조건 \\
\midrule
기준 & C-BAT & 8.50 & 8.93 & 4.37 & -- \\
\midrule
잡음 & 가동유형별 4상태 잡음 & 8.76 & 9.18 & 6.50 & 가동일 편향, 고피크 \\
 & 슬롯 4상태 잡음 & 8.96 & 10.47 & 6.80 & 가동일 편향, 예측 구간 적중률, 고피크 \\
 & Student-$t$ 혁신 & 8.70 & 9.90 & 4.53 & 가동일 편향, 예측 구간 적중률 \\
 & 가동유형별 잡음 + $t$ 혁신 & 8.99 & 10.49 & 7.30 & 네 조건 모두 \\
\midrule
기준일 & 가동유형 기준 & 7.80 & 9.06 & 5.02 & 전체 피크 \\
 & 시간별 계획 기준 & 7.63 & 9.80 & 5.16 & 가동일 편향, 고피크 \\
 & 어텐션 기준일, 전달항 없음 & 7.92 & 11.69 & 4.85 & 전체 피크 \\
 & 어텐션 기준일 + 전달항 & 7.76 & 11.93 & 4.57 & 전체 피크 \\
\midrule
피크 계층 & v1 기준일 피크 회귀 & 7.76 & 11.85 & 21.70 & 전체 피크, 고피크 \\
 & v2 C-BAT 피크 사전 & 7.76 & 12.27 & 20.73 & 전체 피크, 고피크 \\
\midrule
앙상블 & M2 + C-BAT & 8.82 & -- & -- & 슬롯 MAE(C-BAT 8.49) \\
\bottomrule
\end{tabular}
\end{table}
\hwpnote{예측 구간 적중률 조건은 슬롯 90\% 예측 구간 적중률이 0.90에 더 가까워질 것을 요구하였다. 어텐션 기준일은 과거 28일 완전 관측일의 곡선을 시간별 가동 여부와 날짜유형의 유사도로 가중평균한 기준곡선이고, 피크 계층은 C-BAT 곡선과 별도로 일 피크를 회귀로 예측한 2계층 구조이다. 이 실험들은 개발 구간을 다시 쓴 탐색적 결과이다.}

확장 실험은 곡선이나 한 층만 개선하였고 C-BAT의 피크 균형을 넘은 변형은 없었다. 기준일을 바꾼 네 실험은 슬롯 MAE를 7.63--7.92 kW로 낮추었지만 피크 MAE가 9.06--11.93 kW로 커졌다. 어텐션 기준일 두 모형은 가동일 편향(평균 기준)을 C-BAT의 $+9.74$ kW에서 $+14.26$, $+14.74$ kW로 키웠다. 잡음 변형 4종은 예측 구간을 넓혀 경로 최댓값을 더 끌어올렸고, 가동일 편향(중앙값 기준)이 $+9.22$ kW에서 $+9.38$--$+12.45$ kW로 오히려 커졌다. 전용 피크 계층은 비가동일 피크 MAE(v1 1.07 kW)와 C90 경보 F1(v1 0.686)을 개선하였으나, 회귀가 가동일 피크를 평균 쪽으로 수축시켜 고피크일을 20 kW 넘게 낮게 보았다. M2와의 점예측 앙상블은 fold마다 고른 M2 가중치가 0.05--0.30으로 흔들렸고, fold 하나씩 뺀 슬롯 MAE가 8.82 kW로 C-BAT(8.49 kW)보다 나빴다. 시험한 입력과 모형에서는 평범한 가동일의 과대예측을 줄이면서 고피크일 성능을 유지하는 방법을 찾지 못하였다.

이상의 오류분석에서 다음 보완 대상을 정하였다.
\begin{hwplist}
\item 평범한 가동일의 피크 과대예측: 경로 최댓값이 모든 가동일에 약 20 kW의 상승폭을 비슷하게 더하는 것이 주원인이다. 시험한 생산계획 특성과 모형 변형으로는 고피크일과 평범한 가동일을 구분하지 못하였으므로, 이 자료에 없는 설비별 가동 여부·제품 종류·교대조 정보(1.2절)를 확보하여 다시 검증해야 한다.
\item 13--19시간 가동일의 생산 시작 전·종료 뒤 슬롯: C-BAT가 네 모형 중 가장 큰 오차(20.66 kW)를 낸 조건이다. 기준일을 가동유형이나 시간별 계획으로 고르면 곡선은 좋아졌으나 피크가 나빠졌으므로, 곡선 개선과 피크 균형을 함께 만족하는 기준일 규칙이 필요하다.
\item 1--12시간 가동일의 과소예측: 슬롯 편향이 $-5.84$, $-6.14$ kW이고 기존 BAT보다 오차가 컸다. 개발 구간에 6일뿐이므로 표본을 늘려 다시 확인한다.
\item 계절에 따른 피크 수준 변화: 7월의 C90 미탐지와 8월 하순의 과대예측은 모두 피크 수준의 변화를 늦게 따라간 결과로 보인다. 운영 중 최근 관측 피크와 예측 피크의 차이를 감시하고, 학습 자료가 쌓이면 다시 추정한다.
\item 경보 확률의 보정: 7일 Platt 보정이 불안정하여 발행은 보정 전 경로 확률로 하였다. 보정 전 확률은 평범한 가동일에서 높게 나오므로 경보 문턱은 미탐지와 오경보의 비용을 함께 반영하여 정하고, 문턱 결정에 쓰지 않은 기간에서 확인해야 한다.
\item 생산계획 입력: 시각 오차의 영향이 생산량 오차보다 크므로 현장 적용 시 가동 시작·종료 시각을 정확히 입력하고 계획과 실적의 차이를 기록한다.
\end{hwplist}
\hwpnote{조건별 표본이 작고(고피크일 13일, 07--08시 시작 7일, 가동유형 1 6일) 조건별 차이의 신뢰구간은 계산하지 않았다. 개발 구간은 모형 선택에도 쓰였으므로 이 장의 결과는 탐색적이며, 관측된 차이를 인과요인이나 다른 기간의 성능으로 일반화하지 않았다.}
